Photoplethysmography (PPG) has been widely used for drowsiness monitoring, yet most existing studies treat it primarily as a source of conventional physiological features or as input to classification models, leaving the dynamical reorganization of the signal during the Awake--Drowsy transition insufficiently characterized. This study investigates ultra-short PPG dynamics using complementary nonlinear time-series analysis methods, including Simplex Projection, Recurrence Quantification Analysis (RQA), and the Largest Lyapunov Exponent (LLE). Twenty recording sessions from 10 volunteers were analyzed, with 60-s non-overlapping windows used as the primary configuration and 30, 120, and 180~s windows used for robustness analysis. Pseudoperiodic surrogate testing showed that both Awake and Drowsy PPG contained temporal organization beyond that reproduced by a noisy pseudoperiodic null model, without implying deterministic chaos. In the primary 60-s analysis, Drowsiness was associated with reduced forecastability ($\Delta CC=-0.0339$, $q_{\mathrm{BH}}=0.0376$; $\Delta NRMSE=+0.0294$, $q_{\mathrm{BH}}=0.0255$), reduced diagonal recurrence organization ($\Delta DET=-0.0186$, $q_{\mathrm{BH}}=0.0376$), and reduced local trajectory divergence ($\Delta LLE=-0.0376~\mathrm{s}^{-1}$, $q_{\mathrm{BH}}=0.0050$). These directions were preserved across 30--180~s windows, although uncertainty increased at 30~s. Overall, the Awake--Drowsy transition was characterized by coordinated changes in several complementary dynamical properties of PPG rather than by a simple shift along a single ``more chaos''--``less chaos'' axis. These findings provide a dynamical characterization of short-window PPG during declining vigilance and support the future development of nonlinear descriptors for wearable drowsiness monitoring.